\clearpage\section{Hamilton--Jacobi-teori}\label{sec:hj}
\textbf{Fra kanoniske transformationer til løsning af bevægelsen.} I kapitel~\ref{sec:canon} fandt vi, hvordan man kan ændre fasevariable uden at ændre den fysiske bane. Kan vi vælge generatoren så godt, at den \emph{nye Hamiltonfunktion er nul}? Hvis $K(Q,P,t)=0$, giver Hamiltons ligninger straks
\[
\dot Q_i=\frac{\partial K}{\partial P_i}=0,\qquad
\dot P_i=-\frac{\partial K}{\partial Q_i}=0.
\]
De nye variable er således konstanter. Hamilton--Jacobi-teorien (HJ) er metoden til at \emph{finde denne transformation} ved at løse en partiel differentialligning for dens generator.

\subsection{Hovedfunktionen og udledningen af HJ-ligningen}\label{sec:hj-equation}
Vælg en tidsafhængig generator af type 2 med egne argumenter $(q,P,t)$:
\[
F_2(q,P,t)=S(q,P,t).
\]
Fra afsnit~\ref{sec:gen-legendre} ved vi, at
\[
p_i=\frac{\partial S}{\partial q_i},\qquad
Q_i=\frac{\partial S}{\partial P_i},\qquad
K=H(q,p,t)+\left.\frac{\partial S}{\partial t}\right|_{q,P}.
\]
Vælg $S$ sådan, at $K=0$. Efter substitution af $p_i=\partial S/\partial q_i$ fås
\key{H\left(q_1,\ldots,q_n,
\frac{\partial S}{\partial q_1},\ldots,\frac{\partial S}{\partial q_n},t\right)
+\frac{\partial S}{\partial t}=0.}
Funktionen $S$ kaldes \emph{Hamiltons hovedfunktion}. Dens gradient giver de gamle impulser, mens dens parameterafledte giver de nye koordinater. Bemærk, at $\partial S/\partial t$ beregnes ved faste $q,P$; den er ikke $\dd S/\dd t$ langs en bane. HJ er en \emph{førsteordens partiel differentialligning} i $(q,t)$, selv om Hamiltons oprindelige bevægelse kan omfatte mange koblede differentialligninger.

Samtidig gælder for den oprindelige fasevirkning langs en løsning
\[
\frac{\dd S}{\dd t}=
\sum_i\frac{\partial S}{\partial q_i}\dot q_i+\frac{\partial S}{\partial t}
=\sum_i p_i\dot q_i-H,
\]
når $P$ holdes konstant. Dermed er hovedfunktionen (op til passende rand- og integrationskonstanter) virkningen evalueret langs den klassiske bevægelse. Dette er den konkrete forbindelse til fasevirkningen fra kapitel~5.

\subsection{Tidsuafhængig Hamiltonfunktion: $S$ versus $W$}
Hvis $H(q,p)$ ikke afhænger eksplicit af $t$, kan man forsøge at separere tiden:
\[
S(q,\alpha,t)=W(q,\alpha)-Et.
\]
Her er $E$ en separationskonstant, der i den autonome opgave er Hamiltonfunktionens konstante værdi; for $n$ frihedsgrader kan de øvrige $n-1$ uafhængige parametre også indgå i $\alpha$. HJ-ligningen bliver
\key{H\left(q,\frac{\partial W}{\partial q}\right)=E.}
$W$ kaldes \emph{Hamiltons karakteristiske funktion}; den er den rumlige del af $S$, ikke en ny Hamiltonfunktion. Ved en frihedsgrad med
$H=p^2/(2m)+V(q)$ løses først
\[
\frac{\partial W}{\partial q}
=p(q,E)=\pm\sqrt{2m[E-V(q)]},
\]
og derpå
\[
\boxed{W(q,E)=\pm\int^q\sqrt{2m[E-V(s)]}\,\dd s.}
\]
Fortegnet vælges for den relevante impulsgren. Integralet er reelt, hvor $E\ge V(q)$; et vendepunkt kan kræve en ny gren. Ved flervariabel separation $W=\sum_iW_i(q_i;\alpha)$ er dette en \emph{ansats}, som skal kontrolleres ved tilbageindsættelse. Den er ikke mulig for alle Hamiltonfunktioner.

\subsection{Et fuldstændigt integral og rekonstruktion af bevægelsen}
Det er ikke tilstrækkeligt at finde én funktion, der opfylder HJ-ligningen: vi skal også udlede \emph{den fysiske bane}. Et \emph{partikulært integral} er en enkelt løsning af HJ, mens et \emph{fuldstændigt integral} $S(q,\alpha,t)$ lokalt indeholder $n$ uafhængige parametre $\alpha_i$, sådan at den blandede afledte matrix er regulær:
\key{\det\left(\frac{\partial^2S}{\partial q_i\partial\alpha_j}\right)\ne0.}
Denne betingelse gør det muligt lokalt at bruge $S$ som en type 2-generator og rekonstruere løsningen via
\[
\boxed{p_i(t)=\frac{\partial S}{\partial q_i},\qquad
\beta_i=Q_i=\frac{\partial S}{\partial\alpha_i},\qquad
P_i=\alpha_i.}
\]
Da $K=0$, er \emph{både} $\alpha_i$ og $\beta_i$ konstante. For givne initialdata $(q(t_0),p(t_0))$ bestemmes de ved de to sæt relationer ved $t_0$; derefter løses $\beta_i=\partial S/\partial\alpha_i$ for $q_i(t)$, og impulserne findes fra $\partial S/\partial q_i$. Regularitetsbetingelsen er lokal; ved kaustikker, vendepunkter eller koordinatsingulariteter må man ofte skifte gren eller koordinatkort.

For $S=W(q,E,\ldots)-Et$ giver afledning efter energiparameteren
\[
\beta_E=\frac{\partial S}{\partial E}
=\left.\frac{\partial W}{\partial E}\right|_q-t,
\quad\text{altså}\quad
\boxed{\frac{\partial W}{\partial E}=t+\beta_E.}
\]
Her skal $q$ og de øvrige parametre holdes faste ved den partielle afledte. Andre parametre giver tilsvarende $\beta_j=\partial W/\partial\alpha_j$. Energien bestemmer ikke automatisk alle øvrige konstanter eller den valgte fortegnsgren.

\begin{method}\textbf{Tidsuafhængig HJ:} (1) kontroller $\partial H/\partial t=0$ og vælg $S=W-Et$; (2) erstat \emph{alle} $p_i$ med $\partial W/\partial q_i$; (3) separér, hvis strukturen tillader det, og integrér for $W$ med de nødvendige parametre; (4) differentier $S$ efter parametrene ved faste $q$; (5) bestem $\alpha_i,\beta_i$ fra initialdata og løsningens domæne; (6) løs for $q(t)$, beregn $p(t)$ og verificér Hamiltons ligninger.\end{method}

\subsection{Fri partikel: et helt HJ-forløb}
Tag $H(q,p)=p^2/(2m)$. En særlig bekvem \emph{fuldstændig} type 2-generator med parameter $P$ er
\[
S(q,P,t)=qP-\frac{P^2}{2m}t.
\]
Den opfylder
\[
\frac1{2m}\left(\frac{\partial S}{\partial q}\right)^2+
\frac{\partial S}{\partial t}
=\frac{P^2}{2m}-\frac{P^2}{2m}=0,
\qquad
\frac{\partial^2S}{\partial q\partial P}=1\ne0.
\]
De gamle og nye variable er
\[
p=\frac{\partial S}{\partial q}=P,\qquad
Q=\frac{\partial S}{\partial P}=q-\frac Pm t.
\]
Med $P=P_0$ og $Q=Q_0$ konstante finder vi
\[
\boxed{p(t)=P_0,\qquad q(t)=Q_0+\frac{P_0}{m}t.}
\]
Dette demonstrerer hele idéen: HJ-generatoren flytter tidsudviklingen ind i en eksplicit tidsafhængig koordinattransformation. Energi $E=P_0^2/(2m)$ fastlægger kun $|P_0|$, ikke fortegnet; begyndelsesdata afgør retningen.

\subsection{Tidsafhængige Hamiltonfunktioner og separation}
Når $H(q,p,t)$ afhænger eksplicit af tiden, kan man normalt \emph{ikke} sætte $S=W-Et$, men den fulde PDE gælder stadig:
\[
H\left(q,\frac{\partial S}{\partial q},t\right)+\frac{\partial S}{\partial t}=0.
\]
En tidsafhængig separationsansats kan være mulig i særlige modeller, men skal begrundes. Et simpelt kontroleksempel er $H(q,p,t)=p^2/(2m)+f(t)$, hvor
\[
S(q,P,t)=qP-\frac{P^2}{2m}t-\int^t f(\tau)\,\dd\tau
\]
opfylder HJ, giver $p=P$ og $Q=q-Pt/m$. Det rene tidsled i $H$ påvirker generatoren, men ikke de klassiske bevægelsesligninger for $q,p$.
\begin{method}\textbf{Tidsafhængig HJ:} Start fra den fulde PDE, vælg kun en ansats hvis den er konsistent med tidsleddene, find et regulært fuldstændigt integral, bestem $Q_i=\partial S/\partial P_i$, og kontrollér både PDE'en, Hamiltons ligninger og den valgte tids- og parameterg ren.\end{method}

\subsection{HJ efter en tidsafhængig kanonisk transformation: eksamen 2020.1}\label{sec:hj-2020}
Den oprindelige eksamen 2020, opgave 1, begynder med den dæmpede oscillator
\[
\mathcal L(q,\dot q,t)=\frac12e^{\mu t/m}(m\dot q^2-kq^2),
\qquad m,\mu,k>0.
\]
En \emph{første} tidsafhængig generator $F_2(q,P,t)=e^{\mu t/(2m)}qP$ giver
$Q=e^{\mu t/(2m)}q$ og den tidsuafhængige Hamiltonfunktion
\[
\widetilde H(Q,P)=\frac{P^2}{2m}+\frac{kQ^2}{2}+
\frac{\mu}{2m}QP.
\]
Der ønskes nu en \emph{anden} kanonisk transformation med $K=0$, fundet fra HJ. Vi arbejder derfor med de mellemliggende \emph{nye} koordinater $(Q,P)$; HJ skal differentieres med hensyn til $Q$, ikke den oprindelige $q$.
Vælg $S_{\mathrm{HJ}}(Q,E,t)=W(Q,E)-Et$. Sæt
$P=\partial W/\partial Q$ og brug energiligningen
\[
\frac1{2m}\left(\frac{\partial W}{\partial Q}\right)^2+
\frac{\mu Q}{2m}\frac{\partial W}{\partial Q}
+\frac{kQ^2}{2}=E.
\]
Gang med $2m$ og fuldend kvadratet:
\[
\left(W_Q+\frac{\mu Q}{2}\right)^2
=2mE-\left(mk-\frac{\mu^2}{4}\right)Q^2.
\]
For $k>\mu^2/(4m)$ defineres $\omega^2=k/m-\mu^2/(4m^2)>0$, så
\[
W_Q=-\frac{\mu Q}{2}
\pm\sqrt{2mE-m^2\omega^2Q^2},
\]
og dermed
\key{W(Q,E)=-\frac{\mu Q^2}{4}
\pm\int^Q\sqrt{2mE-m^2\omega^2s^2}\,\dd s.}
Denne udledning er indholdet af eksamen 2020.1d; udtrykket er lokalt gyldigt på en gren, hvor kvadratroden er reel.

I opgave 1e bestemmes bevægelsen fra \emph{parameterafledningen}, ikke fra $W$ alene:
\[
\beta=\frac{\partial S_{\mathrm{HJ}}}{\partial E}
=\pm\int^Q\frac{m\,\dd s}{\sqrt{2mE-m^2\omega^2s^2}}-t.
\]
Sæt $A=\sqrt{2E/(m\omega^2)}$. Integralet giver på en valgt gren
\[
t+\beta=\pm\frac1\omega\arcsin\left(\frac QA\right).
\]
Efter valg af fase og korrekt stykvis fortsættelse fås
\[
Q(t)=A\cos(\omega t+\varphi_0),\qquad
\boxed{q(t)=A e^{-\mu t/(2m)}\cos(\omega t+\varphi_0).}
\]
$E$ er her den \emph{konstante værdi af den transformerede} $\widetilde H$, ikke uden videre den oprindelige oscillators bevarede mekaniske energi. Startposition og starthastighed bestemmer $A,\varphi_0$ på den fysiske gren. Tilbage-transformationen er nødvendig; uden den mangler dæmpningen.
\begin{method}Ved en forudgående transformation: (1) brug den transformerede Hamiltonfunktion og \emph{dens} koordinater i HJ; (2) fuldend eventuelt kvadratet, find $W_Q$, integrér; (3) brug $\partial W/\partial E=t+\beta$ til tiden; (4) transformér helt tilbage til de oprindelige variable; (5) verificér den oprindelige bevægelsesligning og initialdata.\end{method}

\subsection{Periodiske bevægelser og aktionsvariable: eksamen 2021.4}
\textbf{Afgrænsning.} Eksamen 2021, opgave 4, bruger \emph{aktions- og vinkelvariable} for periodiske bevægelser, ikke den samme opgave som den direkte HJ-separation i 2020.1. Begrebet hænger sammen med $W$, fordi den klassiske impuls $p=\partial W/\partial q$ kan integreres over en fuld periode. Med eksamenssættets konvention defineres \emph{aktionen}
\[
\boxed{J(E)=\oint p(q,E)\,\dd q,\qquad
\nu=\frac{\partial H(J)}{\partial J}=\frac1T.}
\]
$\nu$ måles i omløb pr.\ tid; den tilhørende vinkelhastighed er $2\pi\nu$. Refleksioner ved hårde vægge behandles stykkevist; HJ-generatoren er ikke ét globalt glat udtryk gennem kollisionerne.

\textbf{2021.4a: fri partikel mellem to vægge.} En masse $m$ bevæger sig mellem $x=0$ og $x=L$ og reflekteres elastisk. På de to grene er $p=\pm\sqrt{2mE}$, og
\[
J=\oint p\,\dd x=2L\sqrt{2mE},\qquad
\boxed{H(J)=\frac{J^2}{8mL^2}}.
\]
Dermed er
\[
\nu=\frac{\partial H}{\partial J}
=\frac{J}{4mL^2}
=\frac{|p|}{2mL}.
\]
Det er netop den inverse tid for frem-og-tilbage-bevægelsen $T=2L/(|p|/m)$.

\textbf{2021.4b: frit fald og elastisk refleksion.} En partikel tabes fra højden $h$ og spejles ved jorden. For $E=mgh$ er den positive gren
$p(x,E)=\sqrt{2m(E-mgx)}$. Derfor
\begin{align*}
J(E)&=2\int_0^{E/(mg)}\sqrt{2m(E-mgx)}\,\dd x\\
&=\frac{4\sqrt{2m}}{3mg}E^{3/2},
\end{align*}
og omskrivning giver
\[
\boxed{H(J)=\left(\frac{3g}{4}\sqrt{\frac m2}\,J\right)^{2/3}.}
\]
Differentiér efter $J$ og indsæt $E=mgh$:
\[
\boxed{\nu=\frac{\partial H}{\partial J}
=\frac12\sqrt{\frac{g}{2h}}.}
\]
Dette stemmer med den direkte fald-og-retur-tid $T=2\sqrt{2h/g}$. Begge eksempler demonstrerer, at et \emph{aktionsintegral} kan bestemme periodetiden uden at løse $x(t)$ eksplicit.
\begin{method}Ved periodiske 1D-problemer: identificér de to grene af $p(q,E)$, beregn $J=\oint p\,\dd q$ med korrekt orientering, isolér $H(J)$ og find $\nu=\partial H/\partial J$. Ved hårde refleksioner opdeles integralet i glatte stykker. Kontrollér ved at beregne den direkte periodetid.\end{method}

\subsection{Samlet metodevalg og forbindelse til kanoniske transformationer}
\begin{method}
\textbf{HJ for tidsuafhængigt $H$:} Sæt $S=W-Et$, løs $H(q,\partial_qW)=E$ og find banen fra $\partial W/\partial E=t+\beta_E$. \textbf{HJ for tidsafhængigt $H$:} Løs hele $H(q,\partial_qS,t)+S_t=0$ og brug et regulært fuldstændigt integral. \textbf{Efter en forudgående transformation:} Brug dens nye koordinater og transformér tilbage til sidst. \textbf{Periodisk bevægelse:} Overvej $J=\oint p\,\dd q$ og $\nu=\partial H/\partial J$. Alle metoder kræver kontrol af grene, integrationskonstanter, regularitet og de oprindelige bevægelsesligninger.
\end{method}
Hamilton--Jacobi er således ikke en konkurrerende fysisk teori: den bruger præcis de kanoniske transformationer fra foregående kapitel til at vælge en beskrivelse, hvor dynamikken ligger i selve transformationen og de nye variable er konstanter.
